\chapter{SEC filings}\label{ch:sec}

\section{Introduction}

A company whose securities are registered with the U.S. Securities and Exchange Commission files its periodic reports -- the annual report on form 10-K, the quarterly report on form 10-Q, and the annual reports of a foreign private issuer on forms 20-F and 40-F -- with its financial statements tagged in Inline XBRL, which the Commission has required of every filer since the phase-in of its rule ended in 2021 \parencite{SEC.2018}. As with ESEF in chapter \ref{ch:esef}, every figure of the primary statements is therefore available as data: the concept it reports, the period, the unit and the value.

The Commission publishes those figures twice: as the filings themselves on EDGAR, and as one JSON file per registrant on \verb|data.sec.gov| that holds the facts already extracted from them \parencite{SEC.2026}. Module \emph{sec} reads the second, just as module \emph{esef} reads the xBRL-JSON that filings.xbrl.org publishes beside the ESEF report packages. Neither module parses Inline XBRL, and neither needs a parser installed. The module builds two tables from the company files, one row per filing and one row per fact, whose shared columns are spelled as module \emph{esef} spells them.

A company file is organised by concept and not by filing: under each taxonomy a concept, under each concept a unit, and under each unit every value the registrant has ever reported for it, each with the accession number of the filing it came from. A filing is reconstructed by grouping on that accession, and the reconstruction is exact. A figure that a later filing restates appears again under the accession of the filing that restated it rather than replacing the first, so each filing goes on saying what it said. Of the twelve thousand eight hundred combinations of concept and period of one registrant, more than half were reported by more than one of its filings, which are the comparative figures of each report standing beside the period the report is about. That is what makes the source valuable beyond the figures themselves: it records what every filing said about every period, and so what a later report changed.

What a company file leaves out decides what the data can be asked. It holds the \emph{undimensioned numeric facts} of a filing and not the whole of it. Every breakdown by segment, product, class or geography is missing, as are the narrative facts and the text blocks of the notes, and every concept the registrant created for its own report; only the concepts of the taxonomies \verb|us-gaap|, \verb|ifrs-full|, \verb|dei| and \verb|srt| are carried. Of a quarterly report of 2026 that states 941 facts, 480 of them without a dimension, the company file carries 321. The tables say so by having no column for a dimension or for text, rather than columns that are always empty.

The company files are downloaded when the pipeline runs and never shipped with the package. Works of the United States Government are not subject to copyright protection in the United States. The Commission asks a client to identify itself in the user agent of its requests, and enforces it: a request without a contact is answered with \verb|403 Forbidden| and a page of HTML that a run which did not check the status would store as the data. The pipeline therefore refuses to build a request without a contact, which is the one point where this module is stricter than module \emph{esef}: filings.xbrl.org asks, and the Commission refuses.

\section{Data pipeline}

\subsection{Landing zone}

There is no landing zone. A company file is a few megabytes at most and is written into the Bronze layer once it is known to be one.

\subsection{Bronze layer}

The Bronze layer keeps the company files as they were published, one per registrant, and beside each a small record of when it was collected and what the server said. A file is asked for with the \verb|Last-Modified| of its last answer as \verb|If-Modified-Since|, so the file of a registrant whose filings have not changed comes back as \verb|304 Not Modified| without a byte. A registrant that has never tagged a filing has no company file, and is recorded as absent rather than as an error.

There are two ways in, and which to use follows from how many registrants a lake wants. A lake interested in a few dozen companies asks for one file each. A lake interested in the whole population takes the archive of all of them, about 1.4 gigabytes, which is what the Commission asks a client wanting many companies to do: ten thousand requests for ten thousand companies is what the archive exists to prevent. The archive is unpacked from memory into the Bronze layer, and where a list of registrants is given only those are written. No \verb|Last-Modified| is recorded for a file that came out of the archive, because the date of the archive is not the date of the file, and a later run asking for that company on its own would otherwise be told it had not changed when it may have.

The Commission limits a client to ten requests a second and answers one that asks faster with \verb|429 Too Many Requests|. Such an answer, and a \verb|503|, is taken to mean ``not now'' rather than ``not at all'': the request is repeated up to five times, waiting as long as the server says or longer. Each company file is collected on its own, and one that fails -- a download cut off, an answer that is not a company file -- is logged with its traceback and counted while the others carry on.

\subsection{Silver layer}

The table of filings has one row per filing a company file names, reconstructed from the facts. A filing is identified by its accession number; its period end is the latest period the filing reports on, which for a periodic report is the period it is about and not one of its comparatives; the fiscal year and period are those of the cover page, which is how a registrant whose financial year ends in January names it. The framework is \verb|US-GAAP| or \verb|IFRS|, taken from the taxonomies the filing tagged against, so that a foreign private issuer reporting under IFRS is recognised as such rather than being taken for a US GAAP filer because of where it filed. By default only the periodic reports on forms 10-K, 10-Q, 20-F and 40-F are kept.

The table of facts has one row per fact of a filing, partitioned by the year the filing was made. That year only moves forward, as a fiscal year does not, since one filing states several of them and a registrant that changes its year end reports two fiscal years or none. A fact has the concept, the first and last day of its period, the unit and the value; there is no column for dimensions, text, decimals or language, each because the company files do not carry it.

Both tables are rebuilt in full from the files the Bronze layer holds. There is no watermark, because there is nothing to resume from: a company file is republished whole whenever any of the registrant's filings changes, so a file that changed changes rows of every year it covers. What a run avoids repeating is the collection, which the records of the Bronze layer handle, not the reading. The files are read on the executors, one company file at a time, so the company file is the unit of parallelism and a run scales with the number of registrants collected. A company file that cannot be read does not fail the run: it comes back as a single row carrying its traceback, which the driver logs and counts, since an exception raised on an executor would fail the whole job and with it every other company.

\subsection{Data quality}

The checks the figures are held to write what they found into a third table of the Silver layer, \emph{issues}, with the machinery shared with module \emph{esef} of chapter \ref{ch:esef}. As there, nothing is dropped and nothing is corrected: the tables are what the Commission published, and a reader who wants to know what is wrong with a figure before using it joins the issues to it. A filing with nothing wrong has no row.

The sixteen checks of the module are mostly arithmetic, the relations a statement states on both sides: total assets against the liabilities and the equity, revenue against the cost of revenue and the gross profit, the income before tax against the tax and the net income, and the three cash flows against the movement in cash. A difference there is a tagging fault rather than a difference of opinion, because the scale of a number is an attribute of its tag and invisible on the page. The remaining checks are structural: a period that ends before it begins, a figure tagged twice with two values, a concept in two units.

One check has no counterpart in module \emph{esef}, and it is the one the company files make possible. Since a company file carries every filing's account of every period, what a year's own report said can be set beside what a later report says about the same year. On NVIDIA's 69 filings the check finds 149 restatements. Share counts and amounts per share are passed over, because a share split restates those by arithmetic rather than by judgement; NVIDIA's two splits accounted for 69 further rows that said nothing about its accounts.

Running the checks over NVIDIA's own filings corrected one of them and explained another apparent fault. A US balance sheet presents mezzanine equity between the liabilities and the equity, and leaving it out of the identity reported NVIDIA's reports for 2016 and 2017 as unbalanced by exactly the 87 and 31 million they hold. And three of its filings appear to have been made before the period they report on had ended. The period end of a filing is the latest period any of its facts reports on, so a fact looking forward -- in these three, an accelerated share repurchase settling in the next quarter -- carries it past the day the filing was made. That is a property of how the period end is derived and not of the filings.

\subsection{Gold layer}

There is no Gold layer. As for ESEF, two report templates read the Silver tables as Parquet and do what a Gold table would do for the question each answers. An \emph{analyst factsheet} covers one registrant over some twenty pages. It differs from its ESEF counterpart in three ways that follow from the source. It can take each year's figures either from that year's own annual report or from the latest report that states the year, the second showing restated figures restated. It brings figures per share onto the share basis of the latest year by a factor read from the filings themselves, since the report of the following year states the share count of the year before again. And it finds the flows of a year by a window of 340 to 380 days, because a registrant on a 52/53-week calendar does not close on the same date every year. A \emph{quality report} reads the table of issues over the whole record. Its rates are taken over the annual reports, the only filings every check was asked of, and it counts by the year a filing was made rather than by fiscal year.

\section{Software architecture}

Module \emph{sec} sits in the package \emph{financialreport} and is organised like module \emph{esef}. The part \emph{common} declares the layout of the data lake and the schemas of the tables; \emph{parse} reads a company file into filings and facts in plain Python, so that it can be tested on its own and run on the executors as it is; \emph{bronze} collects the company files, one at a time or from the archive; \emph{silver} builds the two tables; \emph{quality} holds the checks; and \emph{pipeline} runs collection and tables in one call,
\begin{quote}
\verb|sec.update_sec(spark, "paths.yaml", ["0001045810"],|\\
\verb|               contact="me@example.org")|
\end{quote}
which takes the central index keys of the registrants, padded to ten digits as the tables spell them, and the contact the Commission requires. A registrant is found by its ticker in the Commission's own list of tickers and keys. The checks are a call of their own, \verb|update_silver_sec_quality|, made after the pipeline has run.

A registrant's central index key is not linked to its LEI. Neither the Commission nor GLEIF publishes a mapping from one to the other, and the LEI record of an operating company points at its state register instead: NVIDIA's names Delaware's Division of Corporations and a file number that is not its central index key. The combined tables of chapter \ref{ch:financialreport} therefore carry both keys and derive neither from the other.
